\documentclass[10pt,a4paper]{amsart}
\usepackage[margin=19mm]{geometry}
\usepackage{amsmath,amssymb,mathtools}
\usepackage[scr=rsfs,cal=euler]{mathalfa}
\usepackage{xfrac}
\usepackage[unicode,hidelinks,bookmarks=false]{hyperref}
\newcommand{\ie}{i.e.\ }
\newcommand{\cf}{cf.\ }
\newcommand{\resp}{respectively\ }
\newcommand{\Subsection}{\subsection}
\providecommand{\nobreakdash}{\nobreak\hbox{-}}
\setlength{\emergencystretch}{2em}
\multlinegap0pt
\usepackage{bm}
\usepackage{amssymb}
\usepackage{xcolor}
\newtheorem{lem}{Lemma}
\newtheorem{clm}{Claim}
\newtheorem{linenomath}{Linenomath}
\newcommand{\E}[2][]{\ensuremath{\mathbb{E}_{#1}\left[#2 \right]}}
\newcommand{\Prob}[2][]{\ensuremath{\mathbb{P}_{#1} \left(#2 \right)}}
\newcommand{\Var}[2][]{\operatorname{Var}_{#1}\left(#2\right)}
\newcommand{\eps}{\varepsilon}
\title{Growth scales and uniform integrability of branching processes in varying environments}
\author{Tejas Iyer}
\date{Source: arXiv:2607.05332v1 (CC BY 4.0)}
\begin{document}
\maketitle

\noindent\emph{Source and licence.} Growth scales and uniform integrability of branching processes in varying environments. Version arXiv:2607.05332v1 (CC BY 4.0), distributed by arXiv under the Creative Commons Attribution 4.0 International licence, http://creativecommons.org/licenses/by/4.0/, as recorded in the arXiv licence metadata for this exact version and shown on the abstract page https://arxiv.org/abs/2607.05332v1. Full licence notice: \texttt{licenses/arxiv-2607.05332-CC-BY-4.0.txt}.\par\smallskip

\noindent\textit{Editorial scope.} Counted unit: the article's lem lem:summands-diverge (source0.tex lines 407--416). The article's complete original proof of it is delivered with it (source0.tex lines 457--494, one \texttt{proof} environment of 38 lines, which the article places away from the statement). No mathematics is written, repaired or re-derived: the statement and the proof are the article's own lines, in the article's own notation, and no retained line is substituted or re-flowed.  No further source-local context is needed: every cross-reference of every retained line resolves inside the two delivered spans.  Nothing was omitted from any retained span, so the package carries no omission record.  No retained line contains a citation, so the wrapper carries no bibliography and no bibliography entry.  No retained line contains a figure environment, an image-inclusion command or drawing code, so no image is delivered and the package declares no asset dependency.  Editorial formatting only, in addition: an AMS \texttt{amsart} preamble that restates the article's own declarations and macros used by the retained lines, namely the environments \texttt{lem}, \texttt{clm}, \texttt{linenomath} on separate counters, together with the standard packages those lines need.  The retained environments print in this document as Lemma 1, Clm 1, Linenomath 1 in order of appearance, whereas the article prints the same items with the numbering of its own full document.  The complete native source of the article as distributed in the arXiv e-print for this exact version is bundled unchanged as \texttt{sources/204533/source0.tex}.
\begin{lem} \label{lem:summands-diverge}
    Suppose that $(U_{n,i})_{n\in \mathbb{N}_0, i < n}$ is an array of non-negative, independent random variables and let $\boldsymbol{u} = (u_{i})_{i \in \mathbb{N}_0} \in (0, \infty)^{\mathbb{N}_0}$ be a sequence with $\sum_{i=0}^{\infty} u_i = \infty$. Suppose that for some absolute constant $C > 1$, we have $u_{i} \leq C$ and for all $n \in \mathbb{N}$ we have
    \begin{equation} \label{eq:variance-mean-bound}
        \E{U_{n,i}} = u_i \quad \text{ and } \quad \Var{U_{n,i}} \leq Cu_{i}^2 +  \left(C + \sum_{j=i+1}^{n-2} u_{j}\right)u_{i}. 
    \end{equation}
    Then, for any $r \in \mathbb{N}_0, M > C$, $\eps > 0$ there exists $n =  n(\eps, M, r) > r \in \mathbb{N}$ such that 
    \begin{equation} \label{eq:large-prob-to-be-large}
    \Prob{\sum_{i=r}^{n-1} U_{n,i} \leq M/2} < \eps. 
    \end{equation}
\end{lem}

\begin{proof}[Proof of Lemma~\ref{lem:summands-diverge}]
It suffices to prove the result for $r=0$. Indeed, for fixed $r\in\mathbb{N}_0$, apply the case $r=0$ to the reindexed tail array $(U_{n,i+r})_{n>r,\,i<n-r}$ and the sequence $(u_{i+r})_{i\in \mathbb{N}_0}$. The assumptions are inherited from the original array, and $\sum_{i=r}^{\infty}u_i=\infty$.
We start with the following claim:
\begin{clm} \label{clm:div-sum-ineq}
    For all $M>C$, for any $k \in \mathbb{N}$, there exists an increasing sequence $(n_{i})_{i \in \{0, \ldots, k\}}$ of integers, with $n_{0} = -1$, such that, with
    \[
    s_{j} := 
    \sum_{\ell = n_{j-1}+1}^{n_{j}} u_{\ell}
    \]
    for all $j \in \left\{1, \ldots, k \right\}$ we have 
    \begin{equation} \label{eq:means-dom}
     M \leq s_{j} \quad \text{ and } \quad C + \sum_{\ell = j+1}^{k} s_{\ell} \leq s_{j}.
    \end{equation} 
\end{clm}
Given the truth of the claim, for $\eps > 0$, fix $k$ sufficiently large that $\left(\frac{4(C+2)}{4(C+2)+1}\right)^k \leq \eps$ and choose a sequence $(n_{i})_{i = 0, \ldots k}$ such that~\eqref{eq:means-dom} is satisfied. Then, if we set $N -1 = n_{k}$,
\[
S_{n, j} :=
\sum_{\ell=n_{j-1}+1}^{n_{j}} U_{n, \ell} 
\]
note that $\E{S_{n,j}} = s_{j}$ for $j \in \left\{ 1, \ldots, k\right\}$. Moreover, since each of the summands $U_{n, \ell}$ are independent, by~\eqref{eq:variance-mean-bound} and the fact $n -2 = n_{k}-1$
\begin{align*}
\Var{S_{n,j}} & \leq C\left(\sum_{i=n_{j-1}+1}^{n_{j}} u_{i}^2\right) + \sum_{i=n_{j-1}+1}^{n_{j}} \left(C + \sum_{\ell=i+1}^{n_{k}-1} u_{\ell}\right)u_{i}
\\ & \leq  C\left(\sum_{i=n_{j-1}+1}^{n_{j}} u_{i}^2\right) + \sum_{i=n_{j-1}+1}^{n_{j}} \left(C + \sum_{\ell = j+1}^{k} s_{\ell} + \sum_{\ell=i+1}^{n_{j}} u_{\ell}\right)u_{i}
\\ & \stackrel{\eqref{eq:means-dom}}{\leq} C\left(\sum_{i=n_{j-1}+1}^{n_{j}} u_{i}^2\right) + \sum_{i=n_{j-1}+1}^{n_{j}} \left(s_{j} + \sum_{\ell=i+1}^{n_{j}} u_{\ell}\right)u_{i} 
\\ & \leq C\left(\sum_{i=n_{j-1}+1}^{n_{j}} u_{i}^2\right) + 2s_{j}^2 
\leq (C+2) s_{j}^2. 
 \end{align*}
By  Cantelli's inequality, 
\[
\Prob{S_{n,j} \leq M/2} \stackrel{\eqref{eq:means-dom}}{\leq} \Prob{S_{n,j} \leq \frac{s_{j}}{2}} = \Prob{S_{n,j} - s_{j} \leq - \frac{s_{j}}{2}} \leq \frac{(C+2)s_{j}^2}{(C+2)s_{j}^2 + s_{j}^2/4} = \frac{4(C+2)}{4(C+2) + 1}. 
\]
Thus, by independence of the $S_{n,j}$ for different $j$, and non-negativity of the summands,
\begin{linenomath}
\begin{align*}
\Prob{\sum_{i=0}^{n-1} U_{n,i} \leq M/2} & = \Prob{\sum_{j=1}^{k} S_{n,j} \leq M/2} \\ & \leq \Prob{\forall j\in\left\{1, \ldots, k\right\} \; S_{n,j} \leq M/2} \leq \left(\frac{4(C+2)}{4(C+2)+1}\right)^{k} \leq \eps. 
\end{align*}
\end{linenomath}
\end{proof}

\end{document}
